%----------------------------------------------------------------------------------------
%   PACKAGES AND OTHER DOCUMENT CONFIGURATIONS
%----------------------------------------------------------------------------------------

\documentclass[18pt]{extarticle}

\input{structure.cls} % Include the file specifying the document structure and custom commands

%----------------------------------------------------------------------------------------
%   ASSIGNMENT INFORMATION
%----------------------------------------------------------------------------------------

\title{Εισαγωγή στο Bootloader και στη Διαδικασία Εκκίνησης Λειτουργικού Συστήματος\\
Λειτουργικά Συστήματα (ECE ΓΚ702)} % Title of the assignment

\author{
    \footnotesize Χρήστος Φείδας\\
    \footnotesize \src{fidas@upatras.gr} \and
    \footnotesize Ευάγγελος Λάμπρου\\
    \footnotesize \src{e.lamprou@upnet.gr} \and
    \footnotesize Ζεϊμπέκης Βασίλειος\\
    \footnotesize \src{bzeimp@gmail.com}
} % Author name and email address

\date{University of Patras, Department of Electrical and Computer Engineering} % University, school and/or department name(s) and a date

\usepackage{fancyhdr}
\usepackage{hyperref}

%----------------------------------------------------------------------------------------

\begin{document}

\pagestyle{fancy}
\fancyhf{} % sets both header and footer to nothing
\renewcommand{\headrulewidth}{0pt}
\fancyhead{} % clear all header fields
\fancyfoot{} % clear all footer fields
\fancyfoot[L]{Χρήστος Φείδας, Ευάγγελος Λάμπρου, Ζεϊμπέκης Βασίλειος}
\fancyfoot[R]{\thepage}

\maketitle

%----------------------------------------------------------------------------------------
%   INTRODUCTION
%----------------------------------------------------------------------------------------

\section{Εισαγωγή}

Σε αυτό το εργαστήριο θα μελετήσουμε τα θεμελιώδη στάδια της εκκίνησης (boot process) ενός
λειτουργικού συστήματος σε \textbf{x86 32-bit} αρχιτεκτονική με \textbf{BIOS}.

Ο bootloader είναι ο \textbf{πρώτος κώδικας δικής μας ευθύνης} που εκτελείται μετά το BIOS, σε ένα
περιορισμένο περιβάλλον (Real Mode), και έχει ως βασικό στόχο να:
\begin{itemize}
    \item θέσει το σύστημα σε προβλέψιμη κατάσταση (segments, stack),
    \item φορτώσει τον πυρήνα (kernel) από τον δίσκο στη μνήμη,
    \item προετοιμάσει τη CPU ώστε να εκτελέσει 32-bit κώδικα (Protected Mode),
    \item μεταφέρει τον έλεγχο στον kernel.
\end{itemize}

Μπορείτε να βρείτε το συνοδευτικό βιβλίο του Skyl-OS \href{https://github.com/Billyzeim/Skyl-OS/blob/main/docs/book/main.pdf}{εδώ}.

\subsection{Σημείωση για το αντικείμενο}

Ό,τι περιγράφεται παρακάτω αφορά \textbf{legacy BIOS boot} και \textbf{32-bit x86}. Σύγχρονα συστήματα
χρησιμοποιούν συχνά UEFI, όμως για εκπαιδευτικούς λόγους το BIOS είναι ιδανικό ώστε να δούμε με
καθαρό τρόπο τα βασικά επίπεδα αφαίρεσης και τις αρχιτεκτονικές δεσμεύσεις.

%----------------------------------------------------------------------------------------
%   2.1 REAL MODE
%----------------------------------------------------------------------------------------

\section{Booting in Real Mode}

Για να κατανοήσουμε έναν bootloader πρέπει πρώτα να κατανοήσουμε το \textbf{Real Mode}.

Όταν ο υπολογιστής ενεργοποιείται, ο x86 επεξεργαστής εισέρχεται αυτόματα σε Real Mode. Το Real Mode
διατηρείται για λόγους \textbf{backward compatibility} και σήμερα θεωρείται ιστορικά απαραίτητο αλλά
ακατάλληλο για σύγχρονα OS.

\subsection{16-bit εκτέλεση}

Στο Real Mode, ο επεξεργαστής εκτελεί από προεπιλογή σε \textbf{16-bit} περιβάλλον. Παρότι νεότερες
CPU επιτρέπουν πρόσβαση σε 32-bit registers, για αρχικά/εκπαιδευτικά παραδείγματα αποφεύγουμε να
χτίζουμε σχεδιασμό πάνω σε αυτό, εκτός από την ίδια τη μετάβαση σε Protected Mode.

\subsection{20-bit φυσική διευθυνσιοδότηση: Segment:Offset}

Στο Real Mode, οι διευθύνσεις μνήμης αναφέρονται σε \textbf{φυσικές} διευθύνσεις, μέσω του σχήματος:

\[
PA = Segment \cdot 16 + Offset
\]

όπου:
\begin{itemize}
    \item το \textbf{Segment} προέρχεται από 16-bit segment register,
    \item το \textbf{Offset} προέρχεται από 16-bit address register.
\end{itemize}

Με αυτό το μοντέλο μπορούμε να περιγράψουμε διευθύνσεις μέχρι λίγο κάτω από \textbf{1MiB}.
Επιπλέον, υπάρχουν \textbf{πολλαπλοί τρόποι} να περιγράψουμε την ίδια φυσική διεύθυνση (π.χ. ένα PA
μπορεί να γραφτεί ως διαφορετικά segment/offset ζεύγη).

\subsection{Γιατί το Real Mode είναι ανεπαρκές για σύγχρονα συστήματα}

Εφόσον μιλάμε για φυσικές διευθύνσεις και δεν υπάρχει προστασία, \textbf{καμία διεργασία δεν
απαγορεύεται} από το να προσπελάσει οποιοδήποτε κομμάτι μνήμης. Αυτό καθιστά δύσκολη:
\begin{itemize}
    \item την προστασία κρίσιμων δεδομένων,
    \item τον ορισμό ownership της μνήμης,
    \item την απομόνωση προγραμμάτων.
\end{itemize}

Σύγχρονες προσεγγίσεις λύνουν αυτό το πρόβλημα με \textbf{virtual memory} (π.χ. paging).
Με virtual addresses, το OS χαρτογραφεί κάθε virtual διεύθυνση σε φυσική, εξασφαλίζοντας ότι
κάθε πρόγραμμα “βλέπει” το δικό του χώρο και δεν μπορεί να προσπελάσει άλλον.

\begin{info}[Σημείωση]
Στο βιβλίο του Skyl-OS η κύρια προσέγγιση μνήμης είναι virtual memory μέσω paging (θα καλυφθεί αργότερα),
αλλά στο παρόν κεφάλαιο χρησιμοποιούμε Protected Mode για να προετοιμάσουμε την εκτέλεση 32-bit kernel code.
\end{info}

Η ύπαρξη του Real Mode οδηγεί ιστορικά στη δημιουργία ενός δεύτερου CPU mode: \textbf{Protected Mode},
στο οποίο ανοίγουν δρόμοι για segmentation με privileges (ring levels), και τελικά για πιο ασφαλή
διαχείριση μνήμης.

%----------------------------------------------------------------------------------------
%   2.2 ROLE OF BOOTLOADER / MBR
%----------------------------------------------------------------------------------------

\section{Ο ρόλος του Bootloader και το MBR}

\subsection{BIOS $\rightarrow$ Bootloader}

Το πρώτο πρόγραμμα που αναλαμβάνει τον έλεγχο κατά το boot είναι το BIOS.
Μεταξύ άλλων, το BIOS εντοπίζει συσκευή εκκίνησης και \textbf{διαβάζει τα πρώτα 512 bytes} της
συσκευής (MBR: Master Boot Record).

Αν τα \textbf{τελευταία 2 bytes} του MBR είναι η υπογραφή \textbf{0x55AA}, το BIOS θεωρεί τη συσκευή
bootable και μεταφέρει τον έλεγχο στον Stage-1 bootloader.

\begin{file}[MBR Signature (Snippet)]
\footnotesize \begin{lstlisting}[language={[x86masm]Assembler}, tabsize=2]
times 510 - ($ - $$) db 0
dw 0xAA55
\end{lstlisting}
\end{file}

\subsection{Πού φορτώνεται ο bootloader;}

Το BIOS φορτώνει τον bootloader στη διεύθυνση \src{0x7C00}. Αυτό είναι ιστορική σύμβαση.
Οι bootloader developers το \textbf{υποθέτουν} για συμβατότητα, και η διεύθυνση είναι πρακτική:
είναι κάτω από το 1MiB και αφήνει χώρο χαμηλότερα για δομές όπως interrupt vectors (που θα δούμε σε άλλα labs).

\subsection{Δομή MBR}

Το MBR έχει κλασικά:
\begin{itemize}
    \item 446 bytes boot code (Stage-1),
    \item 64 bytes partition table (4 entries $\times$ 16 bytes),
    \item 2 bytes signature.
\end{itemize}

\begin{figure}[H]
    \centering
    \begin{tabular}{|l|l|l|}
        \hline
        \textbf{Region} & \textbf{Offset (bytes)} & \textbf{Notes} \\
        \hline
        Boot code \& data & 0x000--0x1BD & 446 B (Stage-1) \\
        \hline
        Partition table entry \#1 & 0x1BE--0x1CD & 16 B \\
        \hline
        Partition table entry \#2 & 0x1CE--0x1DD & 16 B \\
        \hline
        Partition table entry \#3 & 0x1DE--0x1ED & 16 B \\
        \hline
        Partition table entry \#4 & 0x1EE--0x1FD & 16 B \\
        \hline
        Boot signature & 0x1FE--0x1FF & 2 B (0x55AA) \\
        \hline
    \end{tabular}
    \caption{Master Boot Record (512 bytes) layout.}
    \label{fig:mbr_layout}
\end{figure}

\subsection{Τι κάνουν τα “advanced” συστήματα;}

Σε πιο πλήρη OS:
\begin{enumerate}
    \item Stage-1 διαβάζει το partition table,
    \item βρίσκει το active/bootable partition,
    \item φορτώνει το Volume Boot Record (VBR) του partition,
    \item το οποίο γνωρίζει για filesystem και οδηγεί σε επόμενο στάδιο bootloading.
\end{enumerate}

\begin{info}[Σημείωση]
Για λόγους απλότητας, στο παρόν κεφάλαιο/εργαστήριο υποθέτουμε ότι ο Stage-1 φορτώνει τον kernel απευθείας
(χωρίς filesystem parsing).
\end{info}

\subsection{Σύντομη εικόνα partition entry}

Κάθε entry περιέχει πληροφορίες όπως: bootable flag, CHS start/end (legacy), type, First LBA και total sectors.
(Τα CHS/LBA θα μας απασχολήσουν στη συνέχεια.)

%----------------------------------------------------------------------------------------
%   2.3 BIOS INTERRUPTS
%----------------------------------------------------------------------------------------

\section{Αξιοποίηση BIOS Interrupts}

Για να φορτώσουμε kernel, πρέπει να μπορούμε να κάνουμε I/O πριν έχουμε δικούς μας drivers.
Το BIOS προσφέρει ένα σύνολο interrupts για βασικές λειτουργίες (disk, screen, keyboard).

\subsection{Γιατί interrupts;}

Τα interrupts λειτουργούν ως \textbf{ελεγχόμενη διεπαφή} προς κρίσιμες λειτουργίες hardware.
Αυτό:
\begin{itemize}
    \item απλοποιεί ανάπτυξη (δεν γράφουμε driver στο bootloader),
    \item αποσυνδέει τον κώδικα από hardware λεπτομέρειες,
    \item επιτρέπει έλεγχο παραμέτρων πριν την εκτέλεση (σημαντικό για ασφάλεια).
\end{itemize}

\subsection{Παράδειγμα: Εκτύπωση χαρακτήρα με INT 0x10}

\begin{file}[BIOS Video: Teletype Output (Snippet)]
\footnotesize \begin{lstlisting}[language={[x86masm]Assembler}, tabsize=2]
mov ah, 0x0E  ; teletype output function
mov al, 'A'   ; ASCII char
mov bh, 0x00  ; page
mov bl, 0x07  ; color
int 0x10
\end{lstlisting}
\end{file}

%----------------------------------------------------------------------------------------
%   2.4 LBA vs CHS
%----------------------------------------------------------------------------------------

\section{LBA vs CHS Addressing}

Πριν φορτώσουμε kernel από δίσκο, πρέπει να κατανοήσουμε το legacy μοντέλο διευθυνσιοδότησης \textbf{CHS}
(Cylinder-Head-Sector), το οποίο “αντικατοπτρίζει” την γεωμετρία των παραδοσιακών σκληρών δίσκων.

\subsection{Γεωμετρική ιδέα: platters, heads, tracks, sectors}

Οι δίσκοι αποτελούνται από \textbf{platters} (κυκλικοί δίσκοι). Σε κάθε επιφάνεια platter γράφει ένας \textbf{head}.
Άρα, για N platters έχουμε 2N heads.

Φανταστείτε νοητούς \textbf{cylinders} με κοινό κέντρο. Η τομή ενός cylinder με μια επιφάνεια platter είναι ένα \textbf{track}.
Κάθε track χωρίζεται σε \textbf{sectors}, και το sector είναι το μικρότερο addressable unit (συνήθως 512 bytes).

\subsection{CHS $\rightarrow$ 512-byte units}

Για να προσπελάσουμε δεδομένα, δίνουμε:
\begin{itemize}
    \item Cylinder,
    \item Head,
    \item Sector (με αρίθμηση που ξεκινά από 1).
\end{itemize}

Κάθε συνδυασμός CHS αντιστοιχεί σε 512-byte block (και πρακτικά αντιστοιχεί σε ένα LBA block).

\subsection{Τυπικά όρια BIOS CHS}

\begin{itemize}
    \item 1024 cylinders (0--1023, 10 bits)
    \item 256 heads (0--255, 8 bits)
    \item 63 sectors (1--63, 6 bits, 1-based)
    \item 512 bytes per sector
\end{itemize}

Το CHS περιορίζεται περίπου στα \textbf{8.4 GiB} και έχει αντικατασταθεί από LBA, όμως το BIOS ακόμη το προσφέρει.
Επιπλέον, σύγχρονοι δίσκοι δεν εκθέτουν πραγματική γεωμετρία και χρησιμοποιούν emulation layer.

\begin{figure}[H]
    \centering
    \includegraphics[width=0.5\textwidth]{images/CHS_Figure.png}
    \caption{CHS addressing layout: cylinders, heads, and sectors.}
    \label{fig:chs_layout}
\end{figure}

%----------------------------------------------------------------------------------------
%   2.5 LOADING THE KERNEL
%----------------------------------------------------------------------------------------

\section{Loading the Kernel}

Ξεκινάμε “ρίχνοντας” λίγο κώδικα. Μην σας αγχώσει: ακολουθεί ανάλυση γραμμή-γραμμή.

\subsection{Απλό start bootloader}

\begin{file}[bootloader.asm (Snippet)]
\footnotesize \begin{lstlisting}[language={[x86masm]Assembler}, tabsize=2]
[ORG 0x7C00]    ; BIOS loads bootloader here
[BITS 16]       ; start in 16-bit Real Mode

start:
    cli

    mov ax, 0x0700
    mov ss, ax
    mov sp, 0x0000

    xor ax, ax
    mov ds, ax
    mov es, ax

    mov ah, 0x02
    mov al, 1
    mov ch, 0
    mov cl, 2
    mov dh, 0
    mov dl, 0x80
    mov bx, 0x1000
    mov es, bx
    xor bx, bx

    int 0x13
\end{lstlisting}
\end{file}

\subsection{Ανάλυση: γιατί cli;}

Το \src{cli} απενεργοποιεί maskable interrupts. Εφόσον \textbf{δεν έχουμε στήσει interrupt handling}
ακόμη, ένα interrupt σε λάθος στιγμή μπορεί να καταστρέψει την εκτέλεση.

\subsection{Ανάλυση: stack σε Real Mode}

Θέτουμε \src{SS=0x0700} και \src{SP=0x0000}. Η φυσική διεύθυνση stack γίνεται:

\[
PA = 0x0700 \cdot 16 + 0x0000 = 0x7000
\]

Το \src{0x7000} επιλέγεται ώστε να είναι \textbf{πάνω από 512 bytes κάτω} από το \src{0x7C00}, όπου βρίσκεται ο bootloader,
ώστε να μειώσουμε πιθανότητα επικάλυψης.

\begin{info}[Σημείωση]
Το bootloader μπορεί “να φαίνεται” ότι δουλεύει χωρίς το πλήρες setup, αλλά χωρίς stack setup και χωρίς μηδενισμό
registers (\src{AX, DS, ES}) το περιβάλλον γίνεται απρόβλεπτο.
\end{info}

\subsection{Ανάγνωση sectors: INT 0x13 (AH=0x02)}

Με την κλήση \src{int 0x13} ζητάμε από το BIOS να φορτώσει δεδομένα από δίσκο σε μνήμη.
Πριν καλέσουμε το interrupt πρέπει να περάσουμε παραμέτρους μέσω registers.

\begin{figure}[H]
    \centering
    \begin{tabular}{|l|l|l|}
        \hline
        \textbf{Register} & \textbf{Purpose} & \textbf{Value Used} \\
        \hline
        AH & BIOS function number & 0x02 (read sectors) \\
        \hline
        AL & Number of sectors to read & 1 \\
        \hline
        CH & Cylinder (part 1) & 0 \\
        \hline
        CL & Sector (bits 0--5) + cyl high bits (6--7) & 2 \\
        \hline
        DH & Head & 0 \\
        \hline
        DL & Drive & 0x80 (first hard drive) \\
        \hline
        ES:BX & Destination segment:offset & 0x1000:0x0000 (= 0x10000) \\
        \hline
    \end{tabular}
    \caption{INT 0x13, AH=0x02 — Disk Read BIOS Call Parameters}
    \label{fig:int13_params}
\end{figure}

Σημεία προσοχής:
\begin{itemize}
    \item Ο τομέας 1 περιέχει bootloader, άρα ξεκινάμε από \textbf{sector 2}.
    \item Τα 2 high bits του cylinder βρίσκονται στα bits 6--7 του \src{CL}.
    \item \src{DL}: 0x00--0x7F floppy, 0x80+ hard drives.
    \item Το \src{ES:BX} είναι ο στόχος στη μνήμη (Real Mode addressing).
\end{itemize}

\subsection{Έλεγχος αποτυχίας: carry flag}

Το BIOS υποδεικνύει error θέτοντας carry flag. Άρα κάνουμε \src{jc} σε error handler.

\begin{file}[disk\_read.asm (Snippet)]
\footnotesize \begin{lstlisting}[language={[x86masm]Assembler}, tabsize=2]
int 0x13
jc disk_error

disk_error:
    mov ah, 0x0E
    mov al, 'E'
    int 0x10
    hlt
\end{lstlisting}
\end{file}

Ο handler τυπώνει ‘E’ με BIOS video teletype και μετά κάνει \src{hlt}.

%----------------------------------------------------------------------------------------
%   2.6 PROTECTED MODE
%----------------------------------------------------------------------------------------

\section{Entering Protected Mode}

Πριν κάνουμε jump στον kernel, αλλάζουμε CPU mode από Real Mode σε \textbf{Protected Mode}.
Τυπικά:
\begin{itemize}
    \item Δεν είναι “απόλυτη” απαίτηση για κάθε kernel,
    \item αλλά στην πράξη \textbf{σύγχρονοι C compilers} παράγουν \textbf{32-bit binaries},
    \item και σε Protected Mode αποκτούμε πρόσβαση σε έως \textbf{4GiB} address space.
\end{itemize}

Για να μπούμε σε Protected Mode, η x86 απαιτεί να έχει περιγραφεί layout μνήμης μέσω \textbf{GDT}
(Global Descriptor Table). Οι segment selectors χρησιμοποιούν τη GDT για να βρουν τα segments.

\begin{info}[Σημείωση]
Στο Skyl-OS αργότερα χρησιμοποιείται paging (virtual memory) ως κύρια διαχείριση μνήμης.
Στο παρόν σημείο, η GDT ορίζεται “απλά” ώστε να ικανοποιηθεί η απαίτηση της CPU, με overlapping code/data segments.
\end{info}

\subsection{Ορισμός απλής GDT}

\begin{file}[gdt.asm (Snippet)]
\footnotesize \begin{lstlisting}[language={[x86masm]Assembler}, tabsize=2]
gdt_start:
    dq 0x0000000000000000 ; null descriptor
    dq 0x00CF9A000000FFFF ; code segment
    dq 0x00CF92000000FFFF ; data segment
gdt_end:

gdt_desc:
    dw gdt_end - gdt_start - 1
    dd gdt_start
\end{lstlisting}
\end{file}

\subsection{Φόρτωση GDT και ενεργοποίηση Protected Mode}

\begin{file}[protected\_mode.asm (Snippet)]
\footnotesize \begin{lstlisting}[language={[x86masm]Assembler}, tabsize=2]
lgdt [gdt_desc]

mov eax, cr0
or eax, 1
mov cr0, eax

jmp 0x08:protected_mode_start
\end{lstlisting}
\end{file}

Μετά την ενεργοποίηση του PE bit στο \src{CR0}, κάνουμε far jump σε code segment selector (0x08) και label που
τρέχει σε 32-bit mode. (Το “πώς/γιατί” του far jump θα αναλυθεί σε επόμενα κεφάλαια.)

\subsection{BITS 32 και jump στον kernel}

\begin{file}[jump\_to\_kernel.asm (Snippet)]
\footnotesize \begin{lstlisting}[language={[x86masm]Assembler}, tabsize=2]
[BITS 32]
protected_mode_start:
    mov ax, 0x10
    mov ds, ax
    mov es, ax
    mov fs, ax
    mov gs, ax
    mov ss, ax

    jmp 0x08:0x10000
\end{lstlisting}
\end{file}

Η διεύθυνση \src{0x10000} είναι εκεί που φορτώσαμε τον kernel νωρίτερα μέσω \src{INT 0x13}.

%----------------------------------------------------------------------------------------
%   DEVELOPMENT ENVIRONMENT (Lab-style)
%----------------------------------------------------------------------------------------

\section{Προετοιμασία περιβάλλοντος ανάπτυξης}

Αν δεν χρησιμοποιείτε Debian-based Linux, προτείνεται virtual machine.

\subsection{Απαραίτητα Πακέτα}

\begin{commandline}
\begin{verbatim}
$ sudo apt-get update
$ sudo apt-get install build-essential nasm qemu-system-x86 git
\end{verbatim}
\end{commandline}

\subsection{Λήψη πηγαίου κώδικα}

Ο κώδικας του Skyl-OS βρίσκεται στο GitHub.

\begin{commandline}
\begin{verbatim}
$ git clone https://github.com/Billyzeim/Skyl-OS.git
$ cd Skyl-OS
\end{verbatim}
\end{commandline}

\begin{info}[Σημείωση]
Αν έχει δοθεί συγκεκριμένο branch για το bootloader lab από τους διδάσκοντες, κάντε checkout σε αυτό.
\end{info}

\subsection{Compilation και εκτέλεση}

\begin{commandline}
\begin{verbatim}
$ make run
\end{verbatim}
\end{commandline}

%----------------------------------------------------------------------------------------
%   EXERCISES (ALL IN ONE PLACE, QUESTIONS ONLY)
%----------------------------------------------------------------------------------------

\section{Ασκήσεις}

\begin{question}
\textbf{Άσκηση 1}\\
Εξηγήστε γιατί ορίζουμε την αρχή του stack στη θέση μνήμης \src{0x7000}. Τι θα γινόταν αν αλλάζαμε αυτή την τιμή σε \src{0x8000};\\
\textbf{Υπόδειξη:} Δείτε Section 5 (Loading the Kernel), υποενότητες για stack setup.
\end{question}

\begin{question}
\textbf{Άσκηση 2}\\
Ο bootloader δεν δουλεύει καθώς ο αριθμός των bytes που φορτώνουμε από τον δίσκο δεν είναι σωστός.
Βρείτε ποιος είναι ο σωστός αριθμός και διορθώστε τον κώδικα ώστε να τρέχει.\\
\textbf{Υπόδειξη:} Συνδέστε “sectors” $\times$ 512 bytes (Section 5) και τη λογική “φόρτωσε τον kernel σε blocks”.
\end{question}

\begin{question}
\textbf{Άσκηση 3}\\
Γιατί χρειαζόμαστε μερικές φορές δεύτερο στάδιο bootloader;\\
\textbf{Υπόδειξη:} Δείτε Section 3 (MBR) για τα όρια 512 bytes, Stage-1, partition table, VBR.
\end{question}

\begin{question}
\textbf{Άσκηση 4}\\
Ποιοί είναι οι περιορισμοί του Real Mode;\\
\textbf{Υπόδειξη:} Δείτε Section 2 (Real Mode) και τη συζήτηση για φυσικές διευθύνσεις/έλλειψη προστασίας.
\end{question}

%----------------------------------------------------------------------------------------
%   DELIVERABLES
%----------------------------------------------------------------------------------------

\section{Παραδοτέα}

Συντάξτε μία αναφορά σε \src{pdf} που θα περιέχει:
\begin{itemize}
    \item απαντήσεις στις θεωρητικές ερωτήσεις (χωρίς μεγάλα dumps κώδικα),
    \item τα σημεία του κώδικα που αλλάξατε και σύντομη αιτιολόγηση,
    \item (προαιρετικά) screenshots/serial output που δείχνουν ότι τρέχει.
\end{itemize}

Υποβάλετε επίσης ένα \src{zip} (π.χ. upAM.zip) με τα αρχεία που τροποποιήσατε.

Καλή επιτυχία!

\section*{Περισσότερες Πηγές}

\begin{itemize}
    \item \href{https://wiki.osdev.org/Boot_Sequence}{OSDev Wiki - BIOS Boot Process}
    \item \href{https://wiki.osdev.org/Protected_mode}{OSDev Wiki - Protected Mode}
    \item \href{https://wiki.osdev.org/MBR_(x86)}{OSDev Wiki - MBR (x86)}
    \item \href{https://wiki.osdev.org/GDT}{OSDev Wiki - GDT}
\end{itemize}

\end{document}
